\documentclass{situationist-preprint}

\title{Cephalopod Consciousness and Distributed Processing}
\msid{MND-23-b}
\msfamily{MND-23}
\msdiscipline{Philosophy of Mind}
\mstype{Companion manuscript}
\mskeywords{cephalopods; octopus; animal consciousness; distributed processing; unity of consciousness}

\begin{document}
\maketitle

\begin{abstract}
Most of the octopus's neurons lie in its arms, which control complex movements locally. Does this distributed organization weaken the case that cephalopods are conscious, or suggest that each arm is a subject? We argue that it does neither. The markers that support the ascription are shown by the whole animal, so the evidence is evidence for an animal-level subject; the arms' local processing contributes to that subject's experience through its results, as much of human processing does; and the partial-unity model of MND-15 accommodates whatever disunity remains. The verdict of MND-23-a stands: cephalopods are probably conscious, the qualification being due to the absence of homology, not to the distribution of processing.
\end{abstract}

\section{The worry}

MND-23-a concludes that cephalopods are probably conscious (Lemma~6.1): the markers of Definition~3.2 converge, but the nervous system that produces them evolved independently of ours, so the inference rests on analogy rather than homology (Principle~P3). About two thirds of an octopus's neurons are in its arms, and an arm severed from the brain can still perform coordinated reaching \citep{Hochner2012}. Two worries arise. The first is that a system so decentralized lacks the integration that consciousness requires, so that the probability of Lemma~6.1 should be lowered. The second is that, if there is consciousness, it may belong to the arms severally rather than to the animal \citep{GodfreySmith2016}.

\section{The animal is the subject}

\begin{proposition}\label{p:animal}
The evidence for cephalopod consciousness is evidence for an animal-level subject.
\end{proposition}

\begin{proof}
The markers on which Lemma~6.1 of MND-23-a rests are exhibited by the animal: it is the octopus that avoids the chamber where it experienced a noxious injection and prefers the one where it received an anesthetic, and the anesthetic's effect on wound-directed behavior is global \citep{Crook2021}. Place learning, exploration and play involve the coordination of all arms with vision and the central brain \citep{Mather2008}. The best explanation of an animal-level pattern, by Principle~P1 of MND-23-a, posits an animal-level state.
\end{proof}

Distribution does not lower the probability either. What the markers require is integration of information sufficient to produce them (marker M7), and Proposition~6.2 of MND-23-a holds that its anatomical location is irrelevant. Human cognition is itself distributed over semi-autonomous processes, as MND-21 holds; much of human motor control is local and reaches experience only through its results. The octopus arm, which elaborates a command issued by the brain and returns filtered information, fits the same pattern at a larger scale. Visual guidance of arm movements, which requires the brain to integrate where an arm is with what the animal sees, has been demonstrated \citep{Gutnick2011}.

\section{Unity}

\begin{objection}[Arm subjects]
An arm processes touch and taste locally and acts on it without central involvement. If the processing of an arm is rich enough, it may sustain experiences that are not the animal's.
\end{objection}

\begin{reply}
Nothing in the evidence requires such experiences, and the evidence for consciousness is evidence about the animal (Proposition~\ref{p:animal}); so the default is a single subject, as MND-23-a holds. If some arm-related experiences are less integrated with one another than human bodily experiences are, the model of MND-15 applies: each may be co-conscious with the animal's central experiences without being co-conscious with every other, which is partial unity within one stream, not a plurality of subjects \citep{CarlsDiamante2017}. The method of MND-15-b counts one stream in such a case. Whether some local processing is unconscious altogether is an open empirical question, and either answer is compatible with the result.
\end{reply}

\section{Conclusion}

The distributed organization of the octopus bears on the structure of its consciousness, not on its presence. The animal is the subject; its experience may be partially unified; and the qualification \enquote{probably} in the verdict of MND-23-a reflects the independent evolution of the cephalopod nervous system, not the location of its neurons.

\begin{thebibliography}{9}\small
\bibitem[Carls-Diamante(2017)]{CarlsDiamante2017} Carls-Diamante, S. (2017). The octopus and the unity of consciousness. \emph{Biology and Philosophy}, 32(6), 1269--1287.
\bibitem[Crook(2021)]{Crook2021} Crook, R.~J. (2021). Behavioral and neurophysiological evidence suggests affective pain experience in octopus. \emph{iScience}, 24(3), 102229.
\bibitem[Godfrey-Smith(2016)]{GodfreySmith2016} Godfrey-Smith, P. (2016). \emph{Other Minds: The Octopus, the Sea, and the Deep Origins of Consciousness}. Farrar, Straus and Giroux.
\bibitem[Gutnick et~al.(2011)]{Gutnick2011} Gutnick, T., Byrne, R.~A., Hochner, B. and Kuba, M. (2011). \emph{Octopus vulgaris} uses visual information to determine the location of its arm. \emph{Current Biology}, 21(6), 460--462.
\bibitem[Hochner(2012)]{Hochner2012} Hochner, B. (2012). An embodied view of octopus neurobiology. \emph{Current Biology}, 22(20), R887--R892.
\bibitem[Mather(2008)]{Mather2008} Mather, J.~A. (2008). Cephalopod consciousness: Behavioural evidence. \emph{Consciousness and Cognition}, 17(1), 37--48.
\end{thebibliography}

\end{document}
